% Nouveau EGA V (cote 25): archive cote 25, site page 153 (PDF page 154).
% typescript. Draft compared: none.
% AI-transcribed by Claude Opus 5.5 (claude-opus-5-5, adaptive thinking, effort high), Message Batches API, batch ega25b (msgbatch_01A7n5YuXubUPNK1zsepi8RG), 2026-11-13; not yet checked by eye.
\scanpage{25/153}
\mspage{\hand{14}}

notions, mis à part \xout{xxx}\inserted[typed]{la} condition dimensionelle = la condition de
$\underline{O}_X$-régularité, le caractère de x se voit \xout{xxx} \unsure{en|eh} fait sur la nature du
point x dans V(u), i.e. sur l'anneau local $\underline{O}_{V(u),x}$ (éventuellement consi-%
déré comme algèbre sur k). Lorsque (dans le cas d'un corps de base k)
x est rationnel sur k, la notion relative à k coincide avec la notion
absolue, comme on voit \inserted{\unsure{précédemment|préalablement}} comme dans le cas où $\underline{E}$ est inversible. Enfin,
lorsque X est lisse sur \over{S}{Y}, et u: $\underline{E}\to\underline{O}_X$ est encore un homomorphisme avec
$\underline{E}$ loc libre de type fini, on définit les notions correspondantes relative-%
ment à S, en regardant fibre par fibre. On remarquera que si x est
non supersingulier relativement à \over{S}{Y}, alors V(u) est plat sur S \xout{relativ} en
x (car u est $\underline{O}_X$-régulier relativement à \over{S}{Y} ...).

\margin{$\big|$} Bien entendu (à signaler tout au moins en r emarque après les
définitions en forme) on pouvait tout aussi bien partir d'une section u
d'un Module loc libre $\underline{E}$ , les définitions précédentes s'appliquent en
interprétant u comme un homomorphisme $\check{\underline{E}}\longrightarrow\underline{O}_X$ .

Lorsque X est lisse sur \over{S}{Y}, u:$\underline{E}\longrightarrow\underline{O}_X$ est donné avec $\underline{E}$ loc libre
de rang m, alors l'ensemble des \xout{points} \inserted[typed]{zéros de u} singuliers rel à S est l'ensembl
sous-jacent à un sous-préschéma fermé $V(u)^{\mathrm{sing}} = V^1$ bien déterminé de
$V(u)=V^o$ , obtenu ainsi: on considère $d^1u$ section de $P^1_{X/\over{S}{Y}}(\check{\underline{E}})$ , sa
restriction à $V^o$ peut alors s'interpréter comme un homomorphisme
\[ \underline{E}\otimes\underline{O}_{V^o} \longrightarrow \underline{\Omega}^1_{X/\over{S}{Y}}\otimes\underline{O}_{V^o} \]
et $V(u)^{\mathrm{sing}}=V^1$ désignera le préschéma des zéros de \xout{\unsure{ce dernier. Rem}}
\xout{Considérant de même $d^2$ section de $P^2_{X/S}(\underline{E})$ puis $d^2$ $V^1$ on}
\xout{trouve un homomorphisme canonique} l'homomorphisme correspondant
\[ \bigwedge^m\underline{E}\otimes\underline{O}_{V^o} \to \underline{\Omega}^m_{X/\over{S}{Y}}\otimes\underline{O}_{V^o} \]
i.e. d'une section bien déterminée de $\underline{\Omega}^m_{X/\over{S}{Y}}\otimes\mathrm{dét}(E)^{-1}\otimes\underline{O}_{V^o}$ .

\margin{153}
